\section{Pasos intermedios (Intermediate Steps): el mecanismo central del razonamiento}

\subsection{Trayectoria histórica: del entrenamiento al prompting}

Denny Zhou trazó el desarrollo de la idea de los pasos intermedios:

\begin{enumerate}
    \item \textbf{2017, Ling et al.}: propusieron en su artículo resolver problemas matemáticos mediante pasos de razonamiento en lenguaje natural (rationale), entrenando modelos seq2seq para generar los pasos intermedios desde cero y llegar así a la respuesta. Denny Zhou lo calificó de «como un viajero en el tiempo».
    \item \textbf{2021, OpenAI GSM8K}: dando continuidad a la idea de 2017, crearon un conjunto de datos matemático con pasos intermedios para hacer fine-tuning de GPT-3, lo que mejoró de manera considerable la capacidad de razonamiento matemático.
    \item \textbf{2021, Google Brain Scratchpad}: descubrieron de forma independiente una idea similar, pero la aplicaron en el dominio de la síntesis de programas (program synthesis) usando pasos intermedios simbólicos.
    \item \textbf{2022, Chain-of-Thought Prompting (Wei et al.)}: evaluaron de manera extensa el efecto de los pasos intermedios en la etapa de prompting, mostrando mejoras notables en prácticamente todas las tareas de PLN.
\end{enumerate}

\begin{importantbox}{Los pasos intermedios son la clave, no el método concreto}
Ya sea mediante training, fine-tuning o prompting, lo que en realidad funciona son los \textbf{pasos intermedios (intermediate steps)} en sí mismos. Cuando se proporcionan ejemplos que incluyen pasos intermedios, el LLM genera respuestas que también los incluyen; esto es una consecuencia natural de que el LLM, como modelo probabilístico, imita el patrón de la entrada.
\end{importantbox}

\subsection{Fundamento teórico: por qué funcionan los pasos intermedios}

Denny Zhou citó un trabajo teórico de 2024 realizado en colaboración con Branden Srouji, de Stanford:

\begin{itemize}
    \item \textbf{Transformer con pasos intermedios}: basta con que la profundidad supere una \textbf{constante} (independiente de la longitud de la entrada) para poder resolver cualquier problema inherentemente serial (inherently serial).
    \item \textbf{Transformer que produce la respuesta directamente}: o bien requiere una profundidad enorme, o bien es imposible de resolver.
\end{itemize}

\begin{warningbox}{Implicación práctica}
Si un modelo no logra resolver cierto problema, una estrategia eficaz es guiarlo para que genere más pasos intermedios. Además, se pueden invocar herramientas externas para apoyar la generación de los pasos intermedios; esta es precisamente una de las ideas centrales del marco de los LLM Agent.
\end{warningbox}

\subsection{Resumen de la sección}

La idea de los pasos intermedios pasó por una evolución que fue del entrenamiento al fine-tuning y de ahí al prompting, pero su esencia se mantuvo invariable: hacer que el modelo derive la respuesta paso a paso en lugar de saltar directamente a ella. En términos teóricos, los pasos intermedios dotan a un Transformer de profundidad finita de la capacidad de resolver cualquier cómputo serial.

